\chapter{Tools and Structured Output}
\label{ch:tools}

\begin{quote}\itshape
The docstring is not documentation. It is part of the prompt.
\end{quote}

\noindent
A tool is a function the model can decide to call. The thing engineers
consistently underestimate is that the function's name, its type hints, its field
descriptions and its docstring are all sent to the model. They are not
commentary --- they are the specification it routes on.

A tool that works perfectly and is described vaguely will be called at the wrong
times, and nothing in your test suite will tell you.

% ============================================================
\section{What the model actually receives}
\label{sec:tool-schema}
\index{tools!schema}
\index{tool@\texttt{@tool}}
\index{docstring}

\begin{python}[caption={A tool, and the schema it generates}, label={lst:tool}]
from langchain.tools import tool


@tool
def get_weather(city: str) -> str:
    """Return the current weather for a city."""
    return "sunny"
\end{python}

\begin{console}
name:        get_weather
description: 'Return the current weather for a city.'
args:        {'city': {'title': 'City', 'type': 'string'}}
\end{console}

The docstring became the description. The type hint became the schema. Both were
sent to the model.

\mermaidfig{tool-schema-to-prompt}{%
  The three things that travel to the model, and what they determine.}{%
  How a Python function becomes part of the prompt.}

For anything non-trivial, declare the schema explicitly so you can describe each
field:

\begin{python}[caption={An explicit schema with descriptions}, label={lst:args-schema}]
from pydantic import BaseModel, Field


class SearchArgs(BaseModel):
    query: str = Field(description="Natural-language query")
    top_k: int = Field(default=5, ge=1, le=50, description="How many results")


@tool(args_schema=SearchArgs)
def search(query: str, top_k: int = 5) -> str:
    """Search the knowledge base."""
    ...
\end{python}

\begin{console}
{'query': {'description': 'Natural-language query', ...},
 'top_k': {'default': 5, 'description': 'How many results',
           'maximum': 50, 'minimum': 1, 'type': 'integer'}}
\end{console}

The constraints travel too. \code{ge=1, le=50} became \code{minimum} and
\code{maximum} in the schema the model reads --- so the bounds are documentation
for the model and validation for you, from one declaration.

\begin{warning}
Write descriptions for the model, not for yourself. ``Gets orders'' tells a
reader what they already guessed and tells the model nothing about \emph{when} to
call it. ``Find a customer's recent orders by customer id. Use this when asked
about order status, delivery or history'' is longer, is not prose you would
enjoy, and measurably changes behaviour.
\end{warning}

% ============================================================
\section{Async tools, and what a sync one costs}
\label{sec:async-tools}
\index{tools!async}
\index{ainvoke@\texttt{ainvoke}}

\begin{python}[caption={The shape to write in a service}, label={lst:async-tool}]
@tool
async def search_orders(customer_id: str, limit: int = 10) -> list[dict]:
    """Find a customer's recent orders by customer id."""
    async with httpx.AsyncClient() as client:
        r = await client.get(f"/orders/{customer_id}", params={"limit": limit})
        return r.json()
\end{python}

A synchronous tool works: the framework runs it in a thread pool when the agent
is invoked asynchronously. It works and it costs. Chapter~\ref{ch:asyncio}
measured the thread-pool escape hatch as adequate and bounded;
Chapter~\ref{ch:serving} measured what happens when it saturates. In a service,
write tools async.

\begin{warning}
The dangerous case is a tool that \emph{looks} async and is not: an
\code{async def} whose body calls a blocking driver. The framework sees a
coroutine and runs it on the loop, where it blocks everything ---
Chapter~\ref{ch:serving} priced that at 44\% of throughput. \code{ruff}'s
\code{ASYNC} rules catch the obvious cases; a blocking call three layers down in
a library they do not.
\end{warning}

% ============================================================
\section{The round trip}
\label{sec:tool-round-trip}
\index{tools!round trip}
\index{ToolMessage@\texttt{ToolMessage}}

\mermaidfig{tool-round-trip}{%
  One turn with two tool calls. Note that both arrive in a single response.}{%
  The tool-calling round trip, including parallel calls.}

Two things worth drawing out.

\textbf{Multiple tool calls arrive together.} A model can ask for several tools
in one response, and they are independent --- so they run concurrently, which is
Chapter~\ref{ch:asyncio}'s \code{gather} doing useful work inside the framework.

\textbf{The loop ends when the model stops asking.} There is no other exit
condition, which is Chapter~\ref{ch:agents}'s warning and the reason a call limit
belongs in every production agent.

\subsection{Errors are messages, not exceptions}
\label{sec:tool-errors}
\index{tools!error handling}

A tool that raises does not necessarily fail the run. The failure can be returned
to the model as a \api{ToolMessage}, which lets it correct itself --- retry with
different arguments, or explain to the user that the lookup failed.

That is a configuration decision, not a given, and it has a sharp edge:

\begin{warning}
Return an error the model can act on, not a stack trace. A Python traceback in a
\api{ToolMessage} is tokens the model cannot use and may quote at your user.
``No customer found with id 123 --- check the id format, which is CUS-nnnn''
produces a correction. \code{KeyError: 'customer\_id'} produces a retry loop.

This is also a security boundary: a traceback leaks internal paths, library
versions and sometimes credentials into a context you may be logging or
streaming to a browser (Chapter~\ref{ch:security}).
\end{warning}

% ============================================================
\section{Too many tools}
\label{sec:tool-count}
\index{tools!selection quality}
\index{measurement!tool selection}

Selection quality degrades as the tool count rises. The commonly cited threshold
is somewhere around fifteen to twenty, it varies by model, and the mechanism is
straightforward: every tool's name, description and schema is in the context on
every call (Chapter~\ref{ch:context}), and the model is choosing from a longer
list of similar-sounding options.

Three structural answers, none of which is ``write better descriptions'':

\begin{center}
\small
\begin{tabularx}{\linewidth}{@{}lX@{}}
\toprule
\textbf{Approach} & \textbf{Trade} \\
\midrule
Tool selector middleware & A model call to narrow the set before the real call.
  Costs a call; saves context and improves selection. \\[3pt]
Subagent partitioning & Each agent has few tools
  (Chapter~\ref{ch:multi-agent}). Costs the topology overhead measured there. \\[3pt]
MCP namespacing & Tools live behind servers and are loaded per task
  (Chapter~\ref{ch:ecosystem}). Costs operational complexity. \\
\bottomrule
\end{tabularx}
\end{center}

\begin{warning}
\textbf{This is the book's fourth experiment and it has not been run.} The
specification: tool-selection accuracy over thirty fixed queries, with five tools
against twenty, and terse descriptions against descriptive ones --- four cells,
twenty runs each, reporting accuracy and spread.

It needs a real model, because selection is exactly the behaviour a fake cannot
simulate. Until it runs, the fifteen-to-twenty figure above is received wisdom
rather than measurement, and this book labels it as such. Appendix~\ref{app:cheatsheet}'s
table for it stays empty.
\end{warning}

% ============================================================
\section{Structured output}
\label{sec:structured-output}
\index{structured output}
\index{ToolStrategy@\texttt{ToolStrategy}}
\index{ProviderStrategy@\texttt{ProviderStrategy}}

When you need a typed object rather than prose, declare the shape:

\begin{python}[caption={Structured output, folded into the agent loop}, label={lst:structured}]
class Ticket(BaseModel):
    priority: str
    category: str
    summary: str


agent = create_agent(model, tools=[...], response_format=Ticket)
result = agent.invoke({"messages": [...]})
result["structured_response"]        # a Ticket instance
\end{python}

Three strategies exist, and they are real classes you can name explicitly:

\begin{center}
\small
\begin{tabularx}{\linewidth}{@{}lX@{}}
\toprule
\textbf{Strategy} & \textbf{Mechanism} \\
\midrule
\api{ProviderStrategy} & The provider's native structured output. Fastest and
  most reliable; takes a \code{strict} flag. Not every provider has it. \\[3pt]
\api{ToolStrategy} & The schema is offered as a tool the model must call. Works
  everywhere; takes \code{handle\_errors} for validation failures. \\[3pt]
\api{AutoStrategy} & Picks native where available, tool-based otherwise. The
  sensible default. \\
\bottomrule
\end{tabularx}
\end{center}

\begin{note}
In 1.0 the schema is required --- there is no prompted-JSON path. Asking a model
to ``respond only with JSON'' and parsing the result was the pre-1.0 pattern, it
failed a few percent of the time in ways that were tedious to handle, and it is
not coming back. If you see it in an example, date the example
(Chapter~\ref{ch:landscape}).
\end{note}

\begin{warning}
Structured output constrains the \emph{shape}, not the \emph{truth}. A
\code{Ticket} with \code{priority="critical"} validates perfectly whether or not
the ticket is critical. Schema validation is not evaluation, and
Chapter~\ref{ch:observability} is where the difference gets measured.
\end{warning}

% ============================================================
\section{What to take away}
\label{sec:ch6-summary}

\begin{itemize}[leftmargin=1.4em]
\item Name, docstring, type hints and field descriptions are all prompt. Write
      them for the model.
\item Pydantic constraints travel into the schema the model reads --- one
      declaration serves as documentation and validation.
\item Write tools async. An \code{async def} wrapping a blocking driver is worse
      than an honest sync tool.
\item Return errors the model can act on. A traceback is unusable tokens and a
      leak.
\item Past roughly twenty tools, selection degrades --- and the answers are
      structural, not stylistic. \textbf{This is unmeasured in this book and
      labelled as such.}
\item Structured output has three strategies; \api{AutoStrategy} is the default.
      Prompted JSON is a pre-1.0 pattern.
\item A schema constrains shape, never truth.
\end{itemize}

\begin{exercisebox}
\textbf{1.} Write a tool, print its generated schema, and check it is what you
would want to read if you were the model.

\textbf{2.} Take a tool with a terse docstring and one with a thorough one, give
the agent both plus eight distractors, and run thirty queries. You have just run
a quarter of experiment 4 --- record the numbers.

\textbf{3.} Make a tool raise. Return the traceback, then return a sentence the
model can act on. Compare what the agent does next.

\textbf{4.} Force \api{ToolStrategy} and then \api{ProviderStrategy} for the same
schema against a provider that supports both, and compare latency and failure
rate over fifty calls.
\end{exercisebox}

\begin{projectbox}
\textbf{Stage 03 --- the tools, and the harness.} Five async tools with explicit
schemas, plus the scriptable fake model everything later depends on.

Two contract items carry the stage. Each tool's generated schema is asserted
against a snapshot --- because that schema is what the model sees, so a change to
it is a behavioural change that should not pass review silently. And the fake
must override \code{bind\_tools}: Chapter~\ref{ch:interrupts} found that no
in-box fake implements it, so none of them can be used inside an agent that has
tools.
\end{projectbox}
